\section{Discussion}
\label{sec:discussion}

\subsection{Interpretation of pilot metrics}

Holdout accuracy of \num{21.4}\% on fourteen classes exceeds uniform guessing (\num{7.1}\%) yet remains far from operational deployment thresholds.
The macro F1 score of \num{0.123} indicates that precision and recall are uneven across meta-actions, which is expected when only five labeled clips per class constrain the \textsc{mlp} head.
Softmax outputs clustered near uniform values (typical top confidence $\approx\num{7}\%$--\num{8}\%) further suggest that the classifier has not learned sharp decision boundaries at the pilot scale.

These outcomes do not invalidate the pipeline; rather, they quantify the sample complexity required before per-person live guidance becomes reliable.
The exploratory analysis confirms ample unused data (\num{5303} clips), so performance gains likely depend on labeled exposure rather than architectural changes alone.

\subsection{Per-person inference versus scene-level baselines}

Assigning labels per \texttt{track\_id} avoids propagating one scene-level prediction to every visible operator.
Session logs demonstrate multiple concurrent tracks with independent label histories.
Nevertheless, crop quality, occlusion, and identity switches remain error sources that were not fully quantified in the absence of frame-level ground truth on mock video.

\subsection{Operational logging and reproducibility}

Versioned session directories bind each detection to model tag (\texttt{all14\_5each}), embedding model identifier, checkpoint name, and timestamp.
This structure supports comparative studies as sampling scales increase.
Configuration fingerprints prevent accidental reuse of mismatched embeddings and checkpoints---a practical requirement when iteration cycles span hours of CPU extraction.

\subsection{Limitations}

Several limitations should be acknowledged openly.
First, experiments used CPU inference; latency figures may decrease on GPU or edge accelerators but were not measured here.
Second, live evaluation lacked per-frame ground truth, precluding supervised session-level F1.
Third, rare \inhard{} classes cap stratified sampling below the nominal one-hundred-clips target.
Fourth, blocked production classes (\emph{Assemble system}, \emph{No action}) may require exclusion or hierarchical modeling to reflect deployment policy.
Fifth, the manuscript reports completed pilot metrics; scaled \texttt{all14\_100each} results remain pending.

\subsection{Future work}

Subsequent iterations should complete scaled training runs, report holdout and live-session metrics in `har-research/04_Analysis_and_Visualization.ipynb`, compare frozen V-JEPA~2 against DINOv2 heads under identical sampling, and evaluate on held-out \inhard{} subjects not seen during stratified sampling.
The production demo integrates per-person inference through the unified VisionOps FastAPI backend (\texttt{vision-ops-backend}, port~8000) with live UI at \texttt{/live} and human review at \texttt{/har-hitl}.
